\section{Comparación funcional}\label{sec:comparacion:funcional}\label{subsec:hallazgos:funcional}

El catálogo tiene tres preguntas sobre funcionalidades: si el plan gratuito incluye \textit{webhooks}, qué API da IP dedicada (una dirección IP de envío que usa un solo cliente) al menor precio y cuánto tiempo se conservan los registros de envío.
Ninguna se responde con la Datasheet ni con PRIME4API, aunque casi todos los datos que piden están publicados en la página de precios o en el panel de cada API.

Para responderlas se lee la matriz de \textit{features} por plan, una tabla que el ciclo de cada Datasheet rellena con las mismas fuentes (Sección~\ref{sec:propuesta:ciclo}).
Cada celda lleva su cita y su peldaño de la escalera de validación, como el contraste documental o la observación de panel, igual que un campo de la Datasheet, y queda vacía cuando ninguna fuente da el dato.
Ocho de sus filas se registran en las ocho APIs, entre ellas los \textit{webhooks} de eventos, la IP dedicada y la retención de actividad.\footnote{Las ocho matrices, con el literal y la fuente de cada celda, están en \texttt{03\_datos/\allowbreak apis/\allowbreak<api>/\allowbreak features.md} del repositorio de datos del trabajo, \url{https://github.com/danvelcam/tfm-pricing-apis-correo}.}

\subsection{\textit{Webhooks} en el plan gratuito}

PI-E01 pregunta si el plan gratuito incluye \textit{webhooks} de eventos, es decir, avisos HTTP que el proveedor envía cuando un correo se entrega, rebota o se abre.
Como en el caso de los 100.000 correos (Sección~\ref{sec:comparacion:caso}), cuenta como \textit{webhook} el aviso HTTP que el propio servicio de correo envía a una URL del cliente, sin otro servicio intermedio.
Es de clase L, la que se responde leyendo la matriz, y las ocho APIs la contestan.

Cinco de las ocho APIs dan \textit{webhooks} en su plan gratuito o en su cuenta de prueba.
Resend Free y MailerSend Free dan un \textit{webhook} cada uno, y la cuenta de prueba de SendGrid da un \textit{Event Webhook}.
Mailgun los incluye en sus cuatro planes, hasta tres URL por tipo de evento, sin recorte en el plan gratuito.
En Postmark la respuesta es «no consta», porque el panel solo los declara en Pro y Platform.
En Brevo la matriz separa dos cosas: los \textit{webhooks} salientes están en los cinco planes, y los de entrada, que reciben en Brevo eventos de otros sistemas, solo en Professional y Enterprise.

Las dos respuestas atípicas son de las APIs que no tienen un plan gratuito con nombre.
SES no tiene \textit{webhooks} en ningún plan, porque sus eventos salen por otros servicios de AWS, como SNS o EventBridge.
En Mailchimp Transactional la cuenta de prueba no los muestra en el panel y la matriz responde «no consta», aunque la documentación da la operación \texttt{webhooks/add} para la cuenta de pago.

\subsection{IP dedicada al menor precio}

PI-E02 pregunta qué API da IP dedicada al menor precio para 200.000 correos al mes.
Combina un dato de la matriz, la IP dedicada, con el precio de enviar 200.000 correos, que la Datasheet da en todas las APIs salvo SES, y la suma se hace a mano.
SES es la más barata, con 44,95 USD: 20,00 de envío, a 0,0001 USD por destinatario, más 24,95 de una IP estándar.
SES vende la IP sin exigir un volumen mínimo de envío.

Tras SES van Mailgun, con 185 USD (plan Foundation de 100.000 correos a 75, con la IP incluida, más 100.000 de sobrecoste a 1,10 USD el millar), y Mailchimp Transactional, con 189,95 USD (ocho bloques de 25.000 correos a 20 USD más 29,95 de la IP).
Resend cuesta 190 USD (plan Scale de 200.000 correos a 160 más 30 de la IP) y SendGrid 199,95 (plan Pro de 100.000 correos a 89,95, con la IP incluida, más 100.000 de sobrecoste a 0,0011).
Mailchimp adelanta a Resend por cinco céntimos.

Brevo tiene la IP dedicada más barata, 251 USD al año, unos 20,92 al mes, pero solo la vende en el plan Professional.
Ese plan arranca en 150.000 correos por 499 USD, y el volumen contratado que cubre 200.000 es el de 250.000, a 599 USD, así que el coste total queda en unos 620 USD al mes.
Es el coste total más alto de las seis APIs que venden la IP con precio publicado, más del triple que Resend.

Dos APIs quedan fuera de la comparación.
Postmark no es elegible, porque exige 300.000 correos al mes para vender la IP.
MailerSend tampoco, porque su página de precios reserva la IP al plan Enterprise sin publicar su precio.

Dos de las tres cifras más bajas llevan salvedad.
La de Mailchimp Transactional no incluye la cuenta Mailchimp de pago que el producto exige, y con los 25 correos por hora de la cuenta de prueba los 200.000 correos tardarían 333 días.
La de SES no sale de la Datasheet, que no tiene campo para el precio por destinatario, y su entorno de pruebas solo permite 200 correos al día.

\subsection{Retención de la actividad}

PI-E03 pregunta cuánto tiempo conserva cada API los registros de envío, y es de clase X: pide un campo que ACM2YAML v0.3 no tiene.
La matriz tiene valor para siete de las ocho APIs.
Postmark guarda 45 días y Resend 30.
SendGrid guarda entre 3 y 7 días según el plan, y MailerSend entre 1 y 30.
Mailgun guarda los registros entre 1 y 30 días y Mailchimp Transactional 30 días en la cuenta de prueba.
En Brevo el periodo lo fija el cliente entre 1 y 24 meses, y es el único caso de las ocho APIs en que no lo pone el proveedor.
En SES ninguna fuente publica la retención, así que la pregunta no tiene respuesta en ningún sitio.

La pregunta sigue en clase X aunque la matriz tenga siete valores, porque la retención es un límite de uso, el máximo de un nivel de uso medible, y no una \textit{feature}.
A diferencia de la cuota mensual, no se gasta con cada envío, y ACM2YAML v0.3 no tiene campo para ese tipo de límite.

\subsection{Lo que la matriz añade}

Con la matriz el catálogo pasa de 25 a 27 preguntas respondibles, porque se suman las dos de clase L.
Las ocho APIs ganan dos respuestas cada una y quedan entre las 20 de Postmark y las 24 de Resend.
Las preguntas comunes a las ocho APIs pasan de 16 a 18, porque se añaden las dos de clase L, PI-E01 y PI-E02.

El resultado deja una asimetría entre los dos tipos de pregunta.
La capacidad se responde con instrumento, ya sea leyendo campos de la Datasheet o ejecutando PRIME4API, mientras que las \textit{features} se responden leyendo la matriz a mano.
En PI-E02 la IP dedicada de cada API está en la matriz y el precio de los 200.000 correos, en la Datasheet salvo en SES, y ninguna herramienta une las dos mitades.
Lo que haría a PI-E02 respondible con instrumento es un vínculo entre la Datasheet del canal API y el iPricing, el metamodelo con el que SPHERE representa los planes y las \textit{features} que un SaaS publica en su canal web \citep{garciafernandez2024pricing4saas}.
Ese vínculo cae fuera del alcance fijado en el Capítulo~\ref{cap:catalogo}, que modela solo el canal API, y por eso la comparación funcional se queda en lectura trazada, celda a celda, sin cálculo automático.
